% !PRINTED 422
\weberpage{422}
%<<CONT
of $F_n[j(\tau), j(\tau)]$ in powers of $q=e^{\pi i \tau}$ with $Cq^{-2N}$
[§\,69, (4)].

If then we develop the factors of the right-hand side of (12) in the same
way, we can determine both $C$ and $N$.

For a factor of the right-hand side of (1) % [sic] The German print cites (1); the preceding sentence and the displayed factor show that (12) is intended.
\[
j(\tau) - j\left(\frac{c+\partial\tau}{a}\right)
\]
we have the following beginnings of the development
\[
\begin{array}{lll}
1. & -q^{-\frac{2\partial}{a}} e^{-\frac{2\pi i c}{a}}, & \partial > a,\\
2. & q^{-2}, & \partial < a,\\
3. & q^{-2}\left(1-e^{-\frac{2\pi i c}{a}}\right), & \partial = a.
\end{array}
\]

Let us first assume that $n$ is not a square; then case 3. does not occur,
and we have
\[
N = \sum_{\partial > a}^{\partial} \frac{\partial}{e}\varphi(e) + \sum_{\partial < a}^{a} \frac{a}{e}\varphi(e),
\]
or, what is the same,
\weq{13}{N = 2\sum_{\partial > a}^{\partial} \frac{\partial}{e}\varphi(e).}

By 1., $C$ is in any case an $n$th root of unity, and since it is at the same
time a rational number, it must be $=\pm 1$.

If then $n$ is a square, $\varphi(\sqrt{n})$ factors of the form 3. occur, and
there results:
\weq{14}{N = 2\sum_{\partial > a}^{\partial} \frac{\partial}{e}\varphi(e) + \varphi(\sqrt{n}).}

Here too $C$ is indeed different from zero, but not equal to $\pm 1$.
In this case we do not need to determine the value of $C$ more closely.
(It is always, as follows from the theory of the division of the circle, a
divisor of $\sqrt{n}$).

If in particular $n$ is a prime, then
\weq{15}{N = 2n.}

One can always dispose of $x$ and $y$ in such a way that, the conditions of
§ 114 being observed, $n = \frac{1}{4}(y^2+mx^2)$ does not become a square.
To this end let $x$ be taken divisible by 4, and $y$ by 2 but not by 4, and
moreover let $y$ have no odd common divisor with $mx^2$. If then
$\frac{1}{4}(y^2+mx^2)$ is a square $M^2$, then
%CONT>>
